% Section 4: The bridge between forecast loss and decision loss
\subsection{The bridge between forecast loss and decision loss}
\label{sec:bridge}

The arguments of Sections~\ref{sec:rl_for_forecasting} and~\ref{sec:forecasting_for_rl}
converge on a single methodological question, namely what loss should govern the
training of the forecaster when the forecast is destined for use inside a
decision. The classical answer is in \citet{GrangerPesaran2000}. The forecast
evaluation problem and the decision evaluation problem are not separable in
general, and the resolution requires going from forecast-error metrics to
decision-induced metrics. The modern answer is in
\citet{ElmachtoubGrigas2022} and \citet{DontiAmosKolter2017}, which make the
Wald (1950) statistical decision theory differentiable and operational at
scale. This subsection lays out the two answers and the obstacles to combining
them with the time-series-foundation-model substrate of
Section~\ref{sec:tsfm}.

The Granger-Pesaran setup begins from the two-state two-action problem of
Section~\ref{sec:forecasting_rl_intro}. Let $y_t \in \{0, 1\}$ index a Bad and
a Good state, $a_t \in \{0, 1\}$ index a No and a Yes action, and
$\{U_{yb}, U_{yg}, U_{nb}, U_{ng}\}$ be the payoffs. The optimal decision rule
is $a_t = \mathbf 1\{\hat\pi_t > \theta_t\}$ with threshold
$\theta_t = c_t / (1 + c_t)$ and cost-benefit ratio
$c_t = (U_{nb} - U_{yb}) / (U_{yg} - U_{ng} + U_{nb} - U_{yb})$. The realised
economic value of the forecast $\hat\pi$ for the binary event over $T$ periods
is
\begin{equation}
V(\hat\pi, \theta) = \bar a +
  T^{-1} \sum_{t=1}^T H_t (y_t - \theta_t)
  \mathbf 1\{\hat\pi_t > \theta_t\},
\label{eq:gp_value}
\end{equation}
where $\bar a$ is invariant across forecasts and $H_t > 0$ is a positive
weight. Granger and Pesaran show that the supreme solution, namely the
forecast that maximises~\eqref{eq:gp_value} for every threshold $\theta_t$
simultaneously, is $\hat\pi_t = \pi_t$, the true conditional probability of the
event. The supreme solution is the natural target for an independent
forecaster who does not know the downstream decision maker's cost ratio,
because it serves all decision makers equally well.

Two scoring devices arise as soon as one evaluates a forecast at finite $T$.
The Brier score \citep{Brier1950} is the cross-sectional analogue of mean
squared error, $B = T^{-1} \sum_t (y_t - \hat\pi_t)^2$, and admits the
\citet{Murphy1973} decomposition
\begin{equation}
B = (\bar y - \bar y^2)
  + K^{-1} \sum_{k=1}^K T_k (\pi_k - \bar y_k)^2
  - K^{-1} \sum_{k=1}^K T_k (\bar y_k - \bar y)^2,
\label{eq:brier_decomp}
\end{equation}
namely the sum of uncertainty, reliability, and resolution, where the
$K$ bins index the forecast probability and $\bar y_k$ is the within-bin
event frequency. The Kuipers score
$\mathrm{KS}(\theta) = H(\theta) - F(\theta)$ is the hit rate minus the false
alarm rate at threshold $\theta$ and is the directional analogue used in
meteorology. \citet{GrangerPesaran2000} establish the exact algebraic
relationship between the Kuipers score and the
\citet{PesaranTimmermann1992} market-timing test statistic. In the standard
case where the per-period weight $H_t = H$ and the threshold $\theta_t = \theta$
are constant, the average economic value relates to the Kuipers score via
\begin{equation}
V(\hat\pi, \theta) = H \bigl[ (1 - \theta) \bar y\, H(\theta)
  - \theta\, (1 - \bar y)\, F(\theta) \bigr],
\label{eq:gp_ks_link}
\end{equation}
and this expression reduces to a constant multiple of the Kuipers score only
in the knife-edge case $\theta = \bar y$, namely when the threshold coincides
with the unconditional event frequency. Outside this case, a forecast that
ranks higher on Brier or on Kuipers can rank lower in realised economic value,
and two forecasts with identical statistical accuracy can be ordered in
opposite directions by the relevant decision loss.

The statistical complement of the Granger-Pesaran argument is the proper
scoring rule machinery of \citet{GneitingRaftery2007}. Their main
characterisation theorem states that every regular proper scoring rule $S$ on
a convex class of distributions $\mathcal P$ arises from a convex generating
function $G : \mathcal P \to \mathbb R$ via
\begin{equation}
S(P, \omega) = G(P) + G^\star(P, \omega) - \int G^\star(P, \cdot) \,\mathrm dP,
\label{eq:proper_score_rep}
\end{equation}
where $G^\star(P, \cdot)$ is a subtangent of $G$ at $P$, and the rule is
strictly proper if and only if $G$ is strictly convex. The Brier score
corresponds to $G(p) = \|p\|_2^2$ on the simplex, the logarithmic score to
negative Shannon entropy, and the continuous ranked probability score and the
energy score follow from analogous constructions on richer sample spaces. The
divergence associated with $S$, namely
$d(P, Q) = S(Q, Q) - S(P, Q)$, is a Bregman divergence under the convex
function $G$, and it is nonnegative with equality only at $P = Q$ when the rule
is strictly proper.

\citet{Dawid1998} closes the loop between proper scoring rules and decision
theory. Given an outcome space $\Omega$, an action space $\mathcal A$, a utility
function $U(\omega, a)$, and the Bayes act $a_P = \arg\max_a \mathbb E_P[
U(\omega, a)]$ under a forecast distribution $P$, the induced score
$S(P, \omega) = U(\omega, a_P)$ is proper with respect to any convex class
$\mathcal P$ containing $P$. Strict propriety holds when the Bayes act is
unique. This is the construction that connects the forecast loss to the
decision loss. The decision-induced score depends on the entire utility
structure of the downstream problem, while a generic strictly proper rule such
as CRPS does not, and the gap between the two is what
\citet{ElmachtoubGrigas2022} and \citet{DontiAmosKolter2017} address with the
SPO$^+$ surrogate and the implicit-KKT differentiation routine respectively.
A reading of the latter two papers as the Wald (1950) statistical decision
theory made differentiable is the most concise way to summarise the modern
synthesis. Wald defined the loss of a decision rule as the expected payoff
under the true distribution, with the optimal procedure being the minimax rule
against the worst admissible distribution. SPO$^+$ and task-based learning
parametrise the predictor, define a differentiable estimate of the
decision-induced loss, and minimise it by stochastic gradient descent, with the
finite-sample guarantees of \citet{LiuGrigas2021} and \citet{HuangGupta2024}
turning the qualitative recipe into quantitative regret bounds at sample
complexity $n^{-1/4}$ on polyhedra and $n^{-1/2}$ on strongly convex level
sets.

Three obstacles separate the current state from the immediate replacement of
standalone forecasting by decision-focused training across the chapter's
domains of interest. First, the pretraining objective of every existing time
series foundation model is action-unconditional. Chronos
\citep{Ansari2024}, TimesFM \citep{Das2024}, and Lag-Llama \citep{Rasul2023}
all model $p(y_{T+1:T+H} \mid y_{1:T})$, possibly with exogenous covariates, but
never $p(y_{T+1:T+H} \mid y_{1:T}, a_T)$ for a decision-maker's action $a_T$.
The forecast distribution conditions on observed history, not on policy. An
RL agent that uses a TSFM as its world model therefore lacks the architectural
hook through which decision-focused training could backpropagate to the
forecaster. Second, the joint-calibration failure documented by
\citet{PerezDiaz2025} implies that path-dependent decisions, including
threshold-crossing inventory rules and worst-case execution costs, see less
calibrated forecasts than marginal ones. Even when the marginal-CRPS objective
is met, the energy-score-based joint-calibration condition typically fails for
trajectory ensembles drawn from current TSFMs, and the failure manifests
exactly on the path-dependent functionals that economic decisions depend on.

Third, the Lucas critique applies in its sharpest form. A foundation forecaster
trained on observational data cannot distinguish the causal effect of a policy
intervention from the response to a confounding shock, and adding parameters
to the forecaster does not change this. The credible remedies parallel the
identification toolkit of econometrics. Instrumental variables, in the spirit
of \citet{Chernozhukov2018} and the broader machine-learning-instruments
literature, can restore identification of the action-conditional dynamics
under exogeneity restrictions. Structural causal model assumptions, in the
spirit of \citet{Buesing2019} and \citet{Oberst2019}, can substitute correctly
specified causal mechanisms for the missing dynamics model. Action-conditional
pretraining, namely training the forecaster on data with action interventions
explicitly recorded and varied, can substitute observational pretraining when
such data is available. None of these remedies is currently implemented at
scale on a TSFM, which is the central open problem we return to in
Section~\ref{sec:forecasting_rl_open}.

The methodological consequence for the practitioner is that the choice of loss
depends on what the forecaster knows about the downstream decision. When the
decision is well-defined and differentiable, decision-focused training in the
spirit of \citet{ElmachtoubGrigas2022}, \citet{DontiAmosKolter2017}, or the
perturbation-gradient generalisation of \citet{HuangGupta2024} is the right
choice, and the finite-sample rate is at best $n^{-1/4}$ on polyhedra and
$n^{-1/2}$ on strongly convex level sets. When the decision is unknown at
training time or shared across many users with different cost structures, a
strictly proper scoring rule in the sense of \citet{GneitingRaftery2007} is
the right choice, because it elicits the supreme solution
$\hat\pi = \pi$ that serves all downstream decisions. When the decision is
known but the forecast is destined for a TSFM-style world model, current
practice is to combine a strictly proper pretraining objective with a
decision-focused fine-tuning pass, and the finite-sample guarantees of this
two-stage recipe are the subject of \citet{Capitaine2025}'s online
DFL programme and remain a research frontier.
